\section{Revised Proof of Theorem 1: Exponential Mean-Bias Amplification}
\label{appendix:proof_theorem1_revised}

\subsection{Setup and Notation}

Let $\{1,\dots,N\}$ be partitioned into clusters
\[
\{1,\dots,N\} = C_1 \sqcup C_2 \sqcup \cdots \sqcup C_M,
\]
and let $c(i)\in\{1,\dots,M\}$ denote the cluster index of token $i$.
For each cluster $C_m$, let $\mu_m \in \mathbb{R}^D$ be its centroid, and write
\[
x_i = \mu_{c(i)} + \varepsilon_i,
\qquad i=1,\dots,N.
\]

We model $\varepsilon_i$ as random residuals around deterministic centroids.
Define the empirical cluster weights
\[
w_m := \frac{|C_m|}{N},
\]
and the \emph{empirical mixture mean}
\[
\mu := \sum_{m=1}^M w_m \mu_m.
\]
We emphasize that $\mu$ is the mean associated with the \emph{current} set of $N$ tokens, and therefore may vary from batch to batch.

We also define the batch mean
\[
\bar{x} := \frac{1}{N}\sum_{i=1}^N x_i.
\]

Let $A=[a_{ij}] \in \mathbb{R}^{N\times N}$ be a row-stochastic matrix,
\[
a_{ij}\ge 0,
\qquad
\sum_{j=1}^N a_{ij}=1,
\]
and define the attention output
\[
y_i := \sum_{j=1}^N a_{ij} x_j.
\]
The output batch mean is
\[
\bar{y} := \frac{1}{N}\sum_{i=1}^N y_i.
\]

For each row $i$, define the in-cluster attention mass
\[
A_{\mathrm{in}}(i) := \sum_{j\in C_{c(i)}} a_{ij}.
\]

We further define:
\[
B^2 := \sum_{m=1}^M w_m \|\mu_m-\mu\|^2
\qquad \text{(between-cluster centroid dispersion),}
\]
\[
\sigma^2 := \frac{1}{N}\sum_{i=1}^N \mathbb{E}\|\varepsilon_i\|^2
\qquad \text{(average residual energy),}
\]
\[
n_{\min} := \min_{m} |C_m|,
\qquad
M_\mu := \max_m \|\mu_m\|,
\qquad
L := \max_{m\neq m'} \|\mu_m-\mu_{m'}\|.
\]

Finally, define the input and output mean-energy ratios by
\[
R_X^2 := \frac{\|\mu\|^2}{\frac{1}{N}\sum_{i=1}^N \mathbb{E}\|x_i\|^2},
\qquad
R_Y^2 := \frac{\|\mathbb{E}\bar{y}\|^2}{\frac{1}{N}\sum_{i=1}^N \mathbb{E}\|y_i\|^2}.
\]

\subsection{Assumptions}

\paragraph{A1 (Centered, independent residuals).}
The residuals $\varepsilon_1,\dots,\varepsilon_N$ are independent random vectors and satisfy
\[
\mathbb{E}[\varepsilon_i] = 0
\qquad \text{for all } i.
\]

\paragraph{A2 (Comparable second moments).}
There exists a constant $\rho \ge 1$ such that
\[
\mathbb{E}\|\varepsilon_i\|^2 \le \rho\,\sigma^2
\qquad \text{for all } i.
\]

\paragraph{A3 (Softmax-compatible cross-cluster suppression).}
For each $i$, the attention mass outside the source cluster is small:
\[
1 - A_{\mathrm{in}}(i) \le \eta_i,
\qquad
\eta := \max_i \eta_i.
\]
In particular, if $A$ is generated by a softmax attention kernel
\[
a_{ij} = \frac{\exp(\ell_{ij}/\tau)}{\sum_{k=1}^N \exp(\ell_{ik}/\tau)},
\]
and if for every $i\in C_m$, every $j\in C_m$, and every $j'\notin C_m$,
\[
\ell_{ij} - \ell_{ij'} \ge \Delta > 0,
\]
then
\[
1 - A_{\mathrm{in}}(i)
\le
\frac{N-|C_m|}{|C_m|} e^{-\Delta/\tau},
\]
hence
\[
\eta \le \max_m \frac{N-|C_m|}{|C_m|} e^{-\Delta/\tau}.
\]

\paragraph{A4 (Approximate in-cluster averaging).}
For each $i\in C_m$, the restriction of row $i$ to cluster $C_m$ is close to uniform on that cluster:
\[
\left\|
\bigl(a_{ij}\bigr)_{j\in C_m}
-
\frac{A_{\mathrm{in}}(i)}{|C_m|}\mathbf{1}
\right\|_1
\le \delta_i.
\]
Define
\[
\bar{\delta}^2 := \frac{1}{N}\sum_{i=1}^N \delta_i^2.
\]

\subsection{Statement of the Revised Theorem}

\begin{theorem}[Exponential Mean-Bias Amplification, expectation version]
\label{thm:exp_mean_bias_amplification_revised}
Under Assumptions A1--A4, let
\[
\alpha := \rho\Bigl(\frac{1}{n_{\min}} + \bar{\delta}^2 + \eta^2\Bigr).
\]
Then the following hold:

\begin{enumerate}
\item \textbf{Input energy decomposition:}
\[
\frac{1}{N}\sum_{i=1}^N \mathbb{E}\|x_i\|^2
=
\|\mu\|^2 + B^2 + \sigma^2,
\]
hence
\[
R_X^2
=
\frac{\|\mu\|^2}{\|\mu\|^2 + B^2 + \sigma^2}.
\]

\item \textbf{Output mean is preserved up to cross-cluster leakage:}
\[
\|\mathbb{E}\bar{y} - \mu\| \le L\eta,
\]
and therefore
\[
\|\mathbb{E}\bar{y}\| \ge \|\mu\| - L\eta.
\]

\item \textbf{Residual energy contracts in expectation:}
\[
\frac{1}{N}\sum_{i=1}^N
\mathbb{E}\Bigl\|
\sum_{j=1}^N a_{ij}\varepsilon_j
\Bigr\|^2
\le
\alpha\,\sigma^2.
\]

\item \textbf{Output ratio lower bound:}
\[
R_Y^2
\ge
\frac{(\|\mu\|-L\eta)^2}{
\|\mu\|^2 + B^2 + \alpha\sigma^2 + 2M_\mu L\eta + L^2\eta^2 }.
\]
In particular, if $\alpha < 1$ and $\eta$ is sufficiently small, then
\[
R_Y^2 > R_X^2.
\]
Thus, when softmax sharply suppresses cross-cluster mass (so that $\eta$ decays exponentially in $\Delta/\tau$) and in-cluster attention behaves as an averaging operator, the mean-energy ratio is amplified.
\end{enumerate}
\end{theorem}

\subsection{Proof}

\paragraph{Step 1: Input mean and input energy.}

By A1,
\[
\mathbb{E}[\bar{x}]
=
\frac{1}{N}\sum_{i=1}^N \mathbb{E}[x_i]
=
\frac{1}{N}\sum_{i=1}^N \mu_{c(i)}
=
\sum_{m=1}^M w_m \mu_m
=
\mu.
\]

Next, expanding $\|x_i\|^2$ gives
\[
\|x_i\|^2
=
\|\mu_{c(i)}\|^2
+ 2\langle \mu_{c(i)}, \varepsilon_i\rangle
+ \|\varepsilon_i\|^2.
\]
Taking expectation and using $\mathbb{E}[\varepsilon_i]=0$,
\[
\mathbb{E}\|x_i\|^2
=
\|\mu_{c(i)}\|^2 + \mathbb{E}\|\varepsilon_i\|^2.
\]
Averaging over $i$,
\[
\frac{1}{N}\sum_{i=1}^N \mathbb{E}\|x_i\|^2
=
\sum_{m=1}^M w_m \|\mu_m\|^2 + \sigma^2.
\]
Using the standard variance decomposition
\[
\sum_{m=1}^M w_m \|\mu_m\|^2
=
\|\mu\|^2 + \sum_{m=1}^M w_m \|\mu_m-\mu\|^2
=
\|\mu\|^2 + B^2,
\]
we obtain
\[
\frac{1}{N}\sum_{i=1}^N \mathbb{E}\|x_i\|^2
=
\|\mu\|^2 + B^2 + \sigma^2.
\]
This proves the expression for $R_X^2$.

\paragraph{Step 2: Exponential suppression of cross-cluster attention.}

Assume $A$ is generated by softmax weights
\[
a_{ij} = \frac{\exp(\ell_{ij}/\tau)}{\sum_k \exp(\ell_{ik}/\tau)},
\]
and fix $i\in C_m$. If
\[
\ell_{ij} - \ell_{ij'} \ge \Delta
\qquad
(j\in C_m,\ j'\notin C_m),
\]
then
\[
\frac{a_{ij'}}{a_{ij}}
=
\exp\!\left(-\frac{\ell_{ij}-\ell_{ij'}}{\tau}\right)
\le
e^{-\Delta/\tau}.
\]
Summing over all $j'\notin C_m$ and comparing to the in-cluster mass yields
\[
1 - A_{\mathrm{in}}(i)
\le
\frac{N-|C_m|}{|C_m|} e^{-\Delta/\tau}.
\]
Hence the worst-case leakage parameter $\eta$ satisfies
\[
\eta \le \max_m \frac{N-|C_m|}{|C_m|} e^{-\Delta/\tau}.
\]

\paragraph{Step 3: Mean preservation up to leakage.}

For each $i\in C_m$, define the centroid-only output
\[
m_i := \sum_{j=1}^N a_{ij}\mu_{c(j)}.
\]
Then
\[
m_i - \mu_m
=
\sum_{j\notin C_m} a_{ij}(\mu_{c(j)}-\mu_m),
\]
because for $j\in C_m$, $\mu_{c(j)}=\mu_m$.
Therefore,
\[
\|m_i-\mu_m\|
\le
\sum_{j\notin C_m} a_{ij}\|\mu_{c(j)}-\mu_m\|
\le
L\sum_{j\notin C_m} a_{ij}
=
L(1-A_{\mathrm{in}}(i))
\le
L\eta.
\]

Now decompose
\[
y_i = m_i + r_i,
\qquad
r_i := \sum_{j=1}^N a_{ij}\varepsilon_j.
\]
By A1,
\[
\mathbb{E}[r_i] = \sum_{j=1}^N a_{ij}\mathbb{E}[\varepsilon_j] = 0,
\]
so
\[
\mathbb{E}[y_i] = m_i.
\]
Averaging over $i$,
\[
\mathbb{E}[\bar{y}] - \mu
=
\frac{1}{N}\sum_{i=1}^N (m_i - \mu_{c(i)}).
\]
Hence, by the triangle inequality,
\[
\|\mathbb{E}\bar{y} - \mu\|
\le
\frac{1}{N}\sum_{i=1}^N \|m_i-\mu_{c(i)}\|
\le
L\eta.
\]
Therefore,
\[
\|\mathbb{E}\bar{y}\| \ge \|\mu\| - L\eta.
\]

\paragraph{Step 4: Residual-energy contraction in expectation.}

Fix $i$. Since the residuals are independent and centered,
\[
\mathbb{E}\langle \varepsilon_j,\varepsilon_k\rangle = 0
\qquad (j\neq k).
\]
Therefore,
\[
\mathbb{E}\|r_i\|^2
=
\mathbb{E}\Bigl\|
\sum_{j=1}^N a_{ij}\varepsilon_j
\Bigr\|^2
=
\sum_{j=1}^N a_{ij}^2 \,\mathbb{E}\|\varepsilon_j\|^2.
\]
By A2,
\[
\mathbb{E}\|r_i\|^2
\le
\rho \sigma^2 \sum_{j=1}^N a_{ij}^2.
\]

We now bound $\sum_j a_{ij}^2$.
Suppose $i\in C_m$, and write the in-cluster restriction as
\[
u^{(i)} := (a_{ij})_{j\in C_m}
=
\frac{A_{\mathrm{in}}(i)}{|C_m|}\mathbf{1} + q^{(i)},
\]
where
\[
\|q^{(i)}\|_1 \le \delta_i,
\qquad
\mathbf{1}^\top q^{(i)} = 0.
\]
Then
\[
\sum_{j\in C_m} a_{ij}^2
=
\|u^{(i)}\|_2^2
=
\left\|
\frac{A_{\mathrm{in}}(i)}{|C_m|}\mathbf{1}
\right\|_2^2
+
\|q^{(i)}\|_2^2,
\]
because the cross term vanishes by $\mathbf{1}^\top q^{(i)}=0$.
Thus
\[
\sum_{j\in C_m} a_{ij}^2
=
\frac{A_{\mathrm{in}}(i)^2}{|C_m|}
+
\|q^{(i)}\|_2^2
\le
\frac{1}{|C_m|}
+
\delta_i^2,
\]
where we used $A_{\mathrm{in}}(i)\le 1$ and $\|q^{(i)}\|_2\le \|q^{(i)}\|_1\le \delta_i$.

For the out-of-cluster part,
\[
\sum_{j\notin C_m} a_{ij}^2
\le
\left(\sum_{j\notin C_m} a_{ij}\right)^2
=
(1-A_{\mathrm{in}}(i))^2
\le
\eta^2.
\]
Therefore,
\[
\sum_{j=1}^N a_{ij}^2
\le
\frac{1}{|C_m|} + \delta_i^2 + \eta^2
\le
\frac{1}{n_{\min}} + \delta_i^2 + \eta^2.
\]
Hence
\[
\mathbb{E}\|r_i\|^2
\le
\rho \sigma^2
\left(
\frac{1}{n_{\min}} + \delta_i^2 + \eta^2
\right).
\]
Averaging over $i$ yields
\[
\frac{1}{N}\sum_{i=1}^N \mathbb{E}\|r_i\|^2
\le
\rho
\left(
\frac{1}{n_{\min}} + \bar{\delta}^2 + \eta^2
\right)\sigma^2
=
\alpha \sigma^2.
\]

\paragraph{Step 5: Output energy bound.}

Since $y_i = m_i + r_i$ and $\mathbb{E}[r_i]=0$,
\[
\mathbb{E}\|y_i\|^2
=
\|m_i\|^2 + \mathbb{E}\|r_i\|^2.
\]
For $i\in C_m$, by Step 3,
\[
\|m_i\|
\le
\|\mu_m\| + \|m_i-\mu_m\|
\le
\|\mu_m\| + L\eta,
\]
so
\[
\|m_i\|^2
\le
\|\mu_m\|^2 + 2\|\mu_m\|L\eta + L^2\eta^2
\le
\|\mu_m\|^2 + 2M_\mu L\eta + L^2\eta^2.
\]
Averaging over $i$,
\[
\frac{1}{N}\sum_{i=1}^N \|m_i\|^2
\le
\sum_{m=1}^M w_m \|\mu_m\|^2 + 2M_\mu L\eta + L^2\eta^2
=
\|\mu\|^2 + B^2 + 2M_\mu L\eta + L^2\eta^2.
\]
Combining this with Step 4 gives
\[
\frac{1}{N}\sum_{i=1}^N \mathbb{E}\|y_i\|^2
\le
\|\mu\|^2 + B^2 + \alpha\sigma^2 + 2M_\mu L\eta + L^2\eta^2.
\]

\paragraph{Step 6: Lower bound for $R_Y^2$.}

Using Step 3 and Step 5,
\[
R_Y^2
=
\frac{\|\mathbb{E}\bar{y}\|^2}{\frac{1}{N}\sum_i \mathbb{E}\|y_i\|^2}
\ge
\frac{(\|\mu\|-L\eta)^2}{
\|\mu\|^2 + B^2 + \alpha\sigma^2 + 2M_\mu L\eta + L^2\eta^2 }.
\]
This proves the stated lower bound.

Finally, when $\alpha<1$ and $\eta$ is sufficiently small, the denominator contains a strictly smaller residual term than the input denominator
\[
\|\mu\|^2 + B^2 + \sigma^2,
\]
while the numerator remains within $O(\eta)$ of $\|\mu\|^2$.
Hence, for sufficiently small $\eta$, one has
\[
R_Y^2 > R_X^2.
\]
This establishes conditional mean-bias amplification.

\qed

